\documentclass[11pt]{article}

\usepackage[margin=1in]{geometry}
\usepackage{amsmath,amssymb,amsthm,mathtools}
\usepackage{booktabs}
\usepackage{microtype}
\usepackage{enumitem}
\usepackage{hyperref}

\hypersetup{colorlinks=true,linkcolor=blue,citecolor=blue,urlcolor=blue}

\newtheorem{theorem}{Theorem}[section]
\newtheorem{proposition}[theorem]{Proposition}
\newtheorem{lemma}[theorem]{Lemma}
\newtheorem{corollary}[theorem]{Corollary}
\newtheorem{example}[theorem]{Example}
\theoremstyle{definition}
\newtheorem{definition}[theorem]{Definition}
\theoremstyle{remark}
\newtheorem{remark}[theorem]{Remark}

\newcommand{\ZZ}{\mathbb Z}
\newcommand{\CC}{\mathbb C}
\newcommand{\one}{\mathbf 1}
\newcommand{\Dd}{\mathcal D}
\newcommand{\HH}{\mathcal H}
\newcommand{\QQ}{\mathcal Q}
\newcommand{\ee}{\mathrm e}
\newcommand{\muM}{\mu}
\newcommand{\Brow}{B_d}
\newcommand{\Kern}{K_d}

\title{\textbf{M\"obius-Weighted Shared-Frequency Clusters\ in the Physical Cross-Divisor Gram}}
\author{Liang Wang\\
Huazhong University of Science and Technology\\
Wuhan 430074, P.R. China}
\date{August 2026}

\begin{document}
\maketitle

\begin{abstract}
TPC-213 showed that the literal transition residual is generated by one common
physical source and that its divisor cross terms are indexed by shared rational
frequencies.  This paper adds the actual M\"obius-log coefficient to that
frequency picture.  For a common reciprocal prime set, common height, and
smooth emitter, we prove an exact dilation covariance: if $h\mid d$ then the
row at denominator $d$ restricted to frequencies with reduced denominator $h$
is exactly the row at $h$.  Consequently the complete-period physical Gram
factors through reduced-denominator clusters, with cluster coefficient
\(C_h=\sum_{h\mid d}\mu(d)\log(d)/d\).  The zero axis vanishes when the prime
shell lies below the Fourier height, and the same reduction commutes with the
four-packet polarization used in the Gate-B route.  Exact rational row data,
followed by numerical logarithm evaluation, gives a weighted physical/direct-sum ratio $0.5963435557$ for
$\{5,7,35\}$ and $1.2119952513$ for $\{3,5,7,105\}$.  Thus the cluster
reduction is a genuine analytic structure, but it has no universal favorable
sign.  No asymptotic saving, prime-shell reassembly, arithmetic $L^2$ credit, or
twin-prime conclusion is claimed.
\end{abstract}

\paragraph{Claim level.}
\texttt{PROVED\_STRUCTURAL\_L1 / MOBIUS\_CLUSTER\_REDUCTION}.\\
The finite sign directions are exact; the displayed ratios are numerical
observations reconstructed from exact rational rows;
the literal V46 asymptotic cluster bound remains open.

\section{Introduction}

The transition-native twin-prime route contains one reciprocal emitter row for
each squarefree divisor.  In the V46 normalization the row has the schematic
form
\begin{equation}
 A_d(r)=\frac{\muM(d)\log d}{d}
 \sum_{q\in\QQ}\sum_{0<|m|\le dq/H}
 \psi\!\left(\frac{Hm}{dq}\right)
 \one_{r\equiv m\overline q\pmod d}.
 \label{eq:v46-row}
\end{equation}
This is the transition-native object frozen in the V46 research artifact
\cite{WangV46}.  TPC-212 isolated its truncated emitter boundary
\cite{WangTPC212}, and TPC-213 supplied the common-source pullback and the
pairwise frequency-intersection Gram \cite{WangTPC213}.
The common-source pullback from TPC-213 showed that the corresponding physical
kernel is not an orthogonal direct sum.  On a complete period, only equal
rational frequencies interact.  The remaining question is whether the
M\"obius-log coefficients make those interactions cancel in a systematic way.

The answer has two parts.  First, there is an exact algebraic collapse.  Strip
the coefficient from \eqref{eq:v46-row} and call the resulting row $B_d$.
Whenever $h\mid d$, the frequency $r=(d/h)a$ has the same amplitude as the
reduced row $B_h(a)$.  All divisors that are multiples of $h$ therefore share
one row, and their coefficients add before squaring.  This is stronger than a
list of pairwise Gram entries: it identifies the correct reduced object.

Second, the coefficient tail is not sign-definite.  The nested family
$\{5,7,35\}$ produces a substantial finite reduction, while the family
$\{3,5,7,105\}$ produces a finite enhancement.  The examples do not model the
full asymptotic problem, but they rule out a universal theorem whose only input
is the existence of shared-frequency coupling.

The contributions are therefore:
\begin{enumerate}[leftmargin=2em]
 \item an exact dilation covariance for the literal reciprocal cutoff;
 \item an exact reduced-denominator factorization of the common-source Gram;
 \item an exact zero-axis and four-packet compatibility statement;
 \item two independently checked finite coefficient ledgers with opposite
       qualitative behavior.
\end{enumerate}

\section{The literal emitter and its reduced frequencies}

Let $\Dd$ be a finite family of squarefree positive integers.  Let $\QQ$ be a
finite set of integers coprime to every member of $\Dd$, let $H\ge 1$, and let
$\psi$ be any function on the finite set of arguments that occur below.  Define
the coefficient-free reciprocal emitter
\begin{equation}
 B_d(r)=\sum_{q\in\QQ}\sum_{0<|m|\le dq/H}
 \psi\!\left(\frac{Hm}{dq}\right)
 \one_{m\overline q\equiv r\pmod d},
 \qquad r\in\ZZ/d\ZZ.
 \label{eq:emitter}
\end{equation}
The V46 coefficient is kept separately as
\begin{equation}
 c_d=\frac{\muM(d)\log d}{d}.
 \label{eq:coefficient}
\end{equation}
For a common physical variable $u$, put
\begin{equation}
 K_d(u)=c_d\sum_{r\bmod d}B_d(r)\ee\!\left(\frac{ru}{d}\right),
 \qquad K(u)=\sum_{d\in\Dd}K_d(u).
 \label{eq:kernel}
\end{equation}

The cutoff is interpreted in the integer sense, so that
$|m|\le\lfloor dq/H\rfloor$.  This convention is important when the dilation
argument changes variables from $m$ to $kn$.

\begin{lemma}[dilation covariance]
\label{lem:dilation}
Suppose $h\mid d$ and write $d=kh$.  For every $a\in\ZZ/h\ZZ$,
\begin{equation}
 B_d(ka)=B_h(a).
 \label{eq:dilation}
\end{equation}
\end{lemma}

\begin{proof}
In the $d$-row, the congruence
$m\overline q\equiv ka\pmod{kh}$ is equivalent to
$m=kn$ and $n\overline q\equiv a\pmod h$.  Indeed, $q$ is a unit modulo
$kh$, so the first congruence forces $k\mid m$.  The range becomes
$|n|\le hq/H$, and the smooth argument is unchanged:
\[
 \frac{H(kn)}{(kh)q}=\frac{Hn}{hq}.
\]
Term-by-term comparison with \eqref{eq:emitter} proves the identity.
\end{proof}

\begin{corollary}[zero axis]
\label{cor:zero}
If every $q\in\QQ$ satisfies $q<H$, then $B_d(0)=0$ for every $d$.  Hence the
reduced denominator $h=1$ contributes no physical Gram energy.
\end{corollary}

\begin{proof}
The congruence for $r=0$ requires $d\mid m$.  But every nonzero admissible
$m$ satisfies $|m|\le dq/H<d$, which is impossible.
\end{proof}

\section{The cluster factorization theorem}

Let
\begin{equation}
 L=\operatorname{lcm}\{d:d\in\Dd\},
 \qquad
 \mathcal H_* =\{h\ge 2:h\mid d\text{ for some }d\in\Dd\}.
 \label{eq:period}
\end{equation}
For $h\in\mathcal H_*$ define the M\"obius-log tail
\begin{equation}
 C_h=\sum_{\substack{d\in\Dd\\h\mid d}}c_d.
 \label{eq:tail}
\end{equation}

\begin{theorem}[reduced-denominator cluster factorization]
\label{thm:cluster}
Assume the hypotheses of Lemma~\ref{lem:dilation}.  On one complete period,
\begin{equation}
 \sum_{u=0}^{L-1}|K(u)|^2
 =L\sum_{h\in\mathcal H_*}|C_h|^2
   \sum_{\substack{a\bmod h\\(a,h)=1}}|B_h(a)|^2.
 \label{eq:cluster}
\end{equation}
The formula remains valid for complex-valued $\psi$ after replacing the inner
square by the Hermitian norm of the row.
\end{theorem}

\begin{proof}
Every nonzero rational frequency in the $d$-row has a unique reduced form
$a/h$ with $h\mid d$ and $(a,h)=1$.  Its representative in $\ZZ/d\ZZ$ is
$r=(d/h)a$.  Lemma~\ref{lem:dilation} gives
\[
 c_dB_d((d/h)a)=c_dB_h(a).
\]
Collecting all divisors that contain the same reduced frequency therefore gives
\[
 K(u)=\sum_{h\in\mathcal H_*}C_h
 \sum_{\substack{a\bmod h\\(a,h)=1}}
 B_h(a)\ee(au/h),
\]
where the zero-frequency term is absent by Corollary~\ref{cor:zero}.  Distinct
reduced fractions have distinct frequencies modulo one.  Summing their products
over $u=0,\ldots,L-1$ gives zero unless the fractions agree, in which case it
gives $L$.  The surviving terms are exactly \eqref{eq:cluster}.
\end{proof}

\begin{corollary}[relation with the pairwise Gram]
\label{cor:pairwise}
The pairwise cross-divisor Gram is not the fundamental object once the
coefficient is restored.  It is the expansion of the cluster sum:
\begin{equation}
 L\sum_{d,e\in\Dd}c_dc_e
 \sum_{\substack{r\bmod d,\ s\bmod e\\r/d\equiv s/e\pmod 1}}
 B_d(r)B_e(s)
 =\text{right-hand side of \eqref{eq:cluster}}.
 \label{eq:pairwise-expansion}
\end{equation}
In particular, the physical period $L$ must be used consistently for every
pair; replacing it by $\operatorname{lcm}(d,e)$ inside a family sum loses the
number of repeated periods.
\end{corollary}

\begin{proof}
Expand the square in \eqref{eq:cluster}, use Lemma~\ref{lem:dilation}, and
apply complete-period orthogonality.  The final sentence is the same identity
with the common period written as a multiple of each pair period.
\end{proof}

\section{Four-packet compatibility}

The Gate-B route uses the complex polarization $a^{(j)}=\beta+i^jw$.
The cluster theorem is linear in the row before taking the Hermitian square, so
the polarization can be applied after the frequency reduction.

\begin{proposition}[polarization commutes with clustering]
\label{prop:polarization}
For any $\beta,w\in\CC$,
\begin{equation}
 \frac14\sum_{j=0}^3 i^j|\beta+i^jw|^2=\beta\overline w.
 \label{eq:polarization}
\end{equation}
If four physical kernels are formed from the rows
$\beta+i^jw$, then each reduced-frequency cluster may be formed first and
\eqref{eq:polarization} applied afterwards.  The resulting polarized scalar is
not positive in general.
\end{proposition}

\begin{proof}
Expand the square.  The sums of $i^j$ and $i^{2j}$ vanish, while
$i^j\overline{i^j}=1$.  The remaining term is $\beta\overline w$.  All maps
used to form a cluster are linear, so the same calculation applies cluster by
cluster.
\end{proof}

\begin{remark}[scope of the zero axis]
The zero-axis statement concerns the additive frequency $r=0$ in the reciprocal
emitter.  It does not delete the multiplicative principal character, a local
Euler exception, or any other mode in a later character decomposition.
\end{remark}

\section{Exact finite certificate}

The release certificate uses
\begin{equation}
 \QQ=\{11,13,17\},\qquad H=40,
 \qquad \psi(t)=\frac{1}{(1+t^2)^2}.
 \label{eq:fixture}
\end{equation}
The profile is smooth and even, and all finite row entries are rational.  The
two divisor families were chosen to test opposite signs rather than to optimize
a favorable outcome.  Energies use the common family period $L$ in both the
physical and direct-sum columns.

\begin{table}[t]
\centering
\caption{Exact-rational row computation followed by numerical logarithm evaluation.}
\label{tab:certificate}
\begin{tabular}{c|c|c|c|c}
\toprule
divisor family $\Dd$ & $L$ & physical cluster energy & direct-sum energy & ratio \\
\midrule
\{5,7,35\} & 35 & 35.0705402044 & 58.8092884914 & 0.5963435557 \\
\{3,5,7,105\} & 105 & 183.3803044582 & 151.3044743922 & 1.2119952513 \\
\bottomrule
\end{tabular}
\end{table}

For the first family, the active coefficient tail at the three nonzero reduced
denominators is approximately
\[
 C_5=-0.2203062093,\qquad C_7=-0.1764057910,\qquad
 C_{35}=0.1015813732.
\]
The physical square sees these tails before squaring.  The direct-sum column
instead retains the three divisor diagonal energies separately.  For the second
family, the quotient structure involving $105$ changes the tail signs enough to
produce a ratio above one.  The sign directions are exact: in the first family,
the only nonzero cross terms pair a negative prime coefficient with the positive
two-prime coefficient; in the second family all four coefficients are negative
and at least one cross Gram is positive.  The decimal ratios are numerical
observations.  This is a finite sign obstruction, not a lower bound
for the actual V46 object.

The producer records exact rational row hashes, exact rational primitive cluster
norms, exact rational Gram entries, and a formal polynomial in the variables
$\log p$.  The independent checker reconstructs those data without importing
the producer module.  Both normal and optimized Python executions pass, and a
separate exact Gaussian-rational check verifies \eqref{eq:polarization}.

\section{Route evaluation and limitations}

Route A is not applicable: this is the analytic prime/twin-prime session rather
than a dynamical spectral family.  The Route-B structural evaluation is
\begin{center}
\begin{tabular}{ll}
\toprule
emitter dilation covariance & \texttt{PROVED\_EXACT} \\
reduced-denominator cluster factor & \texttt{PROVED\_EXACT} \\
zero-axis scope & \texttt{PROVED\_EXACT} \\
four-packet polarization & \texttt{PROVED\_EXACT\_LINEAR\_EXTENSION} \\
finite cancellation/enhancement signs & \texttt{PROVED\_EXACT\_FINITE\_SIGN} \\
finite energy ratios & \texttt{NUMERICAL\_OBSERVATION} \\
universal favorable sign & \texttt{REFUTED\_SCOPED} \\
literal V46 asymptotic cluster bound & \texttt{OPEN} \\
prime-shell reassembly & \texttt{OPEN} \\
arithmetic advance & \texttt{NO} \\
strict $1/400$ endpoint & \texttt{UNPAID} \\
\bottomrule
\end{tabular}
\end{center}

The exact theorem identifies the correct object for the next estimate, but it
does not estimate the coefficient tails uniformly in the transition band
$Y_0<d\le U$.  It also does not perform the prime-only shell reassembly or
attach a source-backed collective BDH theorem to the four polarized packets.
The finite profile is one explicit Schwartz choice, not the proof of an
asymptotic profile law.  Accordingly, the release earns structural route
advance only: no arithmetic $L^2$ credit, fixed-atom credit, strict
$1/400$ payment, or twin-prime conclusion is claimed.

\section{Conclusion}

TPC-214 turns the shared-frequency observation from TPC-213 into an exact
reduced-denominator theorem.  The common-source physical Gram is governed by
M\"obius-log tails $C_h$, not by independent divisor energies.  This is the
reusable structure for the next arithmetic step.  The two finite certificates
also settle the sign question at the right scope: nested coupling can cancel,
but it can also enhance.  The remaining theorem must therefore control these
tails with the literal smooth profile, four-packet signs, zero axis, and prime
shell all retained.

\bibliographystyle{plain}
\bibliography{references}

\end{document}
